\par\noindent\textbf{(1)} 两臂合场E=E1+E2，平均强度I=I1+I2+2Re⟨E1*(t)E2(t+τ)⟩。若两臂源自同一平稳随机场，则相关函数Γ(τ)=⟨E*(t)E(t+τ)⟩，理想等强I1=I2=I0时I(τ)=2I0[1+Reγ(τ)]，γ=Γ(τ)/Γ(0)。
\par\noindent\textbf{(2)} Wiener–Khinchin定理给Γ(τ)=∫\_\{−∞\}∞S(ν)e\textasciicircum{}\{i2πντ\}dν；因此归一复相干度γ(τ)=∫S(ν)e\textasciicircum{}\{i2πντ\}dν/∫S(ν)dν。相干度的模是延迟时干涉可见度的理想上界，而不是源频谱在某一点的值。
\par\noindent\textbf{(3)} 若S(ν)∝exp[−(ν−ν0)²/(2σν²)]，标准Gaussian Fourier积分给γ(τ)=e\textasciicircum{}\{i2πν0τ\}exp(−2π²σν²τ²)。故条纹振幅包络V(τ)=exp(−2π²σν²τ²)，1/e相干时间τc=1/(√2πσν)，谱越宽相干时间越短。
\par\noindent\textbf{(4)} σν→0时γ→e\textasciicircum{}\{i2πν0τ\}，理想单色波可在任意延迟保持模为1。若两臂强度I1≠I2，实测V=2√(I1I2)/(I1+I2)|γ|，即便|γ|=1也可能小于1；偏振与机械振动还会造成额外降低。
